% TOP_PROBLEM: {"id": 20001035, "title": "Exact anatomy of the BCW circular construction", "category_id": 2, "category": {"id": 2, "name": "combinatorics", "display_name": "Combinatorics", "description": "Counting problems, graph theory, discrete structures.", "slug": "combinatorics", "order_index": 2, "created_at": "2026-07-31T15:26:25.671Z"}, "status": "partially_solved", "problem_number": "AIM-COMBINATORICS-0160", "set_id": 14, "statusQualification": "Uses the standard simple loopless BCW formulation; parallel arcs would make the asserted degree–girth bound false."}
% FIRST_POSTED: 2026-10-01
% ENTRY_KIND: statement_only
% UnsolvedMath ID 20001035; source catalogue code AIM-COMBINATORICS-0160
% !TeX program = lualatex
% Definition and sources only; this file is not an LLM solution attempt.
\documentclass[11pt,a4paper]{article}
\usepackage[margin=25mm]{geometry}
\usepackage{fontspec}
\setmainfont{lmroman10-regular.otf}[BoldFont=lmroman10-bold.otf,
 ItalicFont=lmroman10-italic.otf,BoldItalicFont=lmroman10-bolditalic.otf]
\usepackage{amsmath,amssymb,mathtools}
\usepackage{unicode-math}
\setmathfont{latinmodern-math.otf}
\AtBeginDocument{\let\setminus\smallsetminus}
\usepackage{xurl,hyperref}
\hypersetup{colorlinks=true,urlcolor=blue}
\setcounter{secnumdepth}{0}
\setlength{\parindent}{0pt}
\setlength{\parskip}{5pt}
\setlength{\emergencystretch}{3em}
\begin{document}
\section*{Exact anatomy of the BCW circular construction}\label{20001035}
{\small UnsolvedMath ID 20001035 · AIM-COMBINATORICS-0160 · Combinatorics · AIM Workshop Problem Lists}
\subsection{Definitions and mathematical statement}
A finite simple loopless digraph $D$ is $r$-regular if every vertex has
both indegree and outdegree exactly $r$. Opposite arcs are allowed.
Its directed girth is the length of its shortest directed cycle, with
cycle length measured by the number of arcs. For $r\ge1$ a finite
$r$-regular digraph has a directed cycle, so this girth is finite.
For integers $r\ge1$ and $g\ge2$, let
\[
 m(r,g)=\min\{|V(D)|:D\text{ is a simple loopless }r\text{-regular
                         digraph of directed girth }g\}.
\]
The Behzad--Chartrand--Wall conjecture is
\[
                              m(r,g)=r(g-1)+1.
\]
In particular, every such regular digraph should have at least the
specified number of vertices. Both degree conditions hold at each
vertex, rather than merely on average.

The matching candidate construction has vertex set
$\mathbb Z/N\mathbb Z$, where $N=r(g-1)+1$, and arcs
$i\to i+j$ for $j=1,\ldots,r$ modulo $N$. It is $r$-regular and
has girth $g$, showing what sharpness means in the conjecture.
For a short directed cycle the positive increments would have to sum
to a multiple of $N$, which cannot happen in fewer than $g$ steps.
The problem is the universal lower bound for arbitrary digraphs,
without any cyclic symmetry assumption. The simplicity condition
excludes inflating the degree by repeating an arc.
\subsection{Short English statement}
Is r times one less than the girth, plus one, the minimum order of an r-regular directed graph?
\subsection{Sources}
\begin{itemize}
\item[S1] ULAM.AI and the UnsolvedMath Contributors, supplied problems.json, record 20001035 (AIM-COMBINATORICS-0160), AIM Workshop Problem Lists. Catalogue identity and source transcription; CC BY 4.0.
\url{https://huggingface.co/datasets/ulamai/UnsolvedMath}.
\item[S2] American Institute of Mathematics, The Caccetta-Haggkvist conjecture, item 4. Original workshop question underlying AIM-COMBINATORICS-0160.
\url{https://aimath.org/WWN/caccetta/caccetta.pdf}.
\end{itemize}
\end{document}
